\documentclass{article}
\usepackage[charter]{mathdesign}
\usepackage{amsmath}
\newtheorem{example}{Example}
\newtheorem{definition}{Definition}
\newtheorem{proof}{Proof}
\newtheorem{proposition}{Proposition}
\usepackage{booktabs}
\usepackage{titlesec}

% Paragraph spacing
\setlength{\parskip}{7pt}
\setlength{\parindent}{0pt}

% Equation spacing
\setlength{\abovedisplayskip}{10pt}
\setlength{\belowdisplayskip}{10pt}
\setlength{\abovedisplayshortskip}{6pt}
\setlength{\belowdisplayshortskip}{6pt}

% Section spacing
\titlespacing*{\subsection}
{0pt}{0.5cm}{0pt}


%-------------------------------------------------------------------------------------------------
\begin{document}
	
	\title{Structural Aliasing Analysis in Time Series Forecasting using Chronos} % Article title, use manual lines breaks (\\) to beautify the layout
	\author{
		Matteo Boniotti, Gianluca Brignoli and Federico Sabbadini \thanks{This work is licensed under the Creative Commons Attribution 4.0 International License. \newline Copyright for any third-party material included in this work remains with its respective owners and must be respected accordingly. \newline \newline During the preparation of this work, the authors used AI-based tools, including GPT-5.5, GPT-5.4, Gemini 3.1 Pro, and Sonnet 4.6, to support language refinement and improve clarity. All generated or suggested content was carefully reviewed, revised where necessary, and validated by the authors, who assume full responsibility for the final content of this work.}
	}
	\maketitle
	
	%----------------------------------------ABSTRACT-------------------------------------------
	\begin{abstract}
		This document shows the author year style of citations. A lot of things. A lot of things. A lot of things. A lot of things. A lot of things. A lot of things. A lot of things. A lot of things. A lot of things. A lot of things. A lot of things. A lot of things. A lot of things. A lot of things. A lot of things. A lot of things. A lot of things. A lot of things. A lot of things. A lot of things. A lot of things. A lot of things. A lot of things. A lot of things. A lot of things. A lot of things. A lot of things. A lot of things. A lot of things. A lot of things. A lot of things. A lot of things. A lot of things. A lot of things. A lot of things. A lot of things. A lot of things. A lot of things. A lot of things. A lot of things. A lot of things. A lot of things. A lot of things. A lot of things. A lot of things. A lot of things. A lot of things. A lot of things. A lot of things. A lot of things. A lot of things. A lot of things. A lot of things. A lot of things. A lot of things. A lot of things. A lot of things. 
	\end{abstract}
	
	%----------------------------------------INTRO------------------------------------------------
	\newpage
	\section*{Introduction}
		Time-series foundation models have emerged as a powerful paradigm for zero-shot and few-shot temporal forecasting, with architectures such as Chronos~\cite{FastAccurateZeroshot2024} achieving remarkable generalisation across diverse domains.
A key enabler of this scalability is \emph{patch-based tokenisation}: rather than feeding individual scalar observations to a transformer, contiguous windows of samples are grouped into fixed-length patches, dramatically shortening the effective sequence length and reducing the quadratic cost of self-attention.
Despite its computational appeal, this design choice introduces a subtle but consequential inductive bias that has received little formal scrutiny.

The central concern of this work is what we term \emph{structural aliasing}: a representational failure that arises not from undersampling in the classical Nyquist–Shannon sense, but from the geometric interaction between the periodicity of an input signal and the fixed stride of the patch grid.
When the stride $S$ is commensurate with the signal period $T_0$, consecutive patches become identical or near-identical, collapsing the token sequence to a constant or low-diversity representation.
The model then suffers a loss of temporal observability that is invisible to standard sampling-theoretic guarantees — the signal is fully representable, yet the architecture cannot perceive it.

Beyond the signal-theoretic dimension, structural aliasing carries semantic consequences.
In real-world deployment scenarios, time-series signals encode domain concepts such as operational states, anomalies, control conditions, and recurring events.
A model that is blind to certain frequencies may therefore misclassify or entirely miss semantically meaningful patterns.
To make this connection explicit, we introduce an ontological layer that maps signal-level observability failures to domain-level concept distortions, and we evaluate the phenomenon through both behavioural metrics and internal representational probing.

This report is structured as follows.
Section~\ref{sec:one} formalises patch-based tokenisation as a linear operator and derives the conditions under which token identity — and hence observability loss — occurs.
Section~\ref{sec:two} characterises the resulting blind-spot frequencies as harmonics of the patch frequency $f_{\mathrm{patch}} = f_s / S$, and analyses the role of patch size, stride, and overlap.
Section~\ref{sec:point3} introduces the semantic layer through a formal ontology of the chosen application domain.
Section~\ref{sec:point4} enumerates the falsifiable hypotheses that drive our experimental programme.
Sections~\ref{sec:point5} and~\ref{sec:point6} describe the controlled synthetic experiments and the evaluation methodology at both the behavioural and representational levels.
Section~\ref{sec:point7} presents the Bayesian statistical analysis used to quantify uncertainty and compare competing hypotheses.
Section~\ref{sec:point8} investigates a principled mitigation strategy, and Section~\ref{sec:point9} situates the findings within an agentic AI setting, discussing how an autonomous agent could monitor, explain, and act upon structural aliasing risks.

	
	%----------------------------------------POINT1------------------------------------------------
	\newpage
	\section{Formal Definition of Patch-Based tokenization} \label{sec:one}
		 \subsection{Setup and Formal Definition of the Tokenisation Operator}

To rigorously understand the failure modes of patch-based architectures like Chronos-Bolt, we must first formally define the tokenisation process. Let $x[n]$ represent a discrete-time, univariate signal where the index $n \in \mathbb{Z}$ is sampled at a continuous frequency of $f_s$.

Unlike classical autoregressive models that process temporal data one scalar time-step at a time, modern foundation models employ patch-based tokenisation. This technique groups contiguous samples to radically reduce the effective sequence length and improve the computational efficiency of the transformer's attention mechanisms. This is achieved by sliding a window of fixed size $P$ (the patch length) across the signal, moving forward by a step size of $S$ (the stride).

We can extract this sequence of tokens as follows:
\begin{equation}
    \mathbf{p}_k = \bigl(x[kS],\; x[kS+1],\; \ldots,\; x[kS+P-1]\bigr) \in \mathbb{R}^P, \quad k \in \mathbb{N}
\end{equation}

\begin{definition}[Patch Tokenisation Operator]
    We can formalise this extraction as a linear map acting on the bounded signal space. The map $\mathcal{T}: \ell^\infty(\mathbb{Z}) \to (\mathbb{R}^P)^{\mathbb{N}}$, as defined above, is referred to as the \emph{patch tokenisation operator}. When collecting an arbitrary number of $K$ patches from a signal vector $\mathbf{x}$, the operation can be represented as a matrix multiplication:
    \begin{equation}
        \mathbf{P} = \mathcal{T}\{x\} = \mathbf{W} \cdot \mathbf{x}
    \end{equation}
    Here, $\mathbf{W} \in \{0,1\}^{KP \times N}$ is a block-binary selection matrix. Its $k$-th block consists of $P$ rows that act as a mask, picking out the exactly $P$ consecutive signal samples starting at the temporal index $kS$.
\end{definition}

\subsection{The Mechanics of Phase-Locking and Degeneracy}

The fundamental vulnerability of the patch tokenisation operator emerges when the periodic structure of the input signal aligns perfectly with the model's fixed stride grid. We refer to this specific synchronization as ``phase-locking.''

\begin{proposition}[Condition for Token Identity]
    Two consecutive patches, $\mathbf{p}_k$ and $\mathbf{p}_{k+1}$, will be mathematically identical if and only if the signal amplitude remains constant across the stride shift for every element within the observation window:
    \begin{equation}
        x[n] = x[n + S] \quad \forall n \in \{kS,\, kS+1,\, \ldots,\, kS+P-1\}
    \end{equation}
    Globally, if the underlying continuous signal is a pure sinusoid (or any periodic signal) with a fundamental period of $T_0$ samples, this exact alignment occurs whenever the stride $S$ spans an integer multiple of the signal's period:
    \begin{equation}
        \boxed{S = m \cdot T_0, \quad m \in \mathbb{Z}^+}
    \end{equation}
    By substituting $T_0 = f_s / f$, we can express this condition in terms of the signal's continuous frequency $f$ and the sampling frequency $f_s$:
    \begin{equation}
        f = m \cdot \frac{f_s}{S}, \quad m \in \mathbb{Z}^+
    \end{equation}
\end{proposition}

\begin{proof}
    To satisfy $\mathbf{p}_k = \mathbf{p}_{k+1}$ strictly component-wise, it is required that $x[kS + j] = x[(k+1)S + j]$ for every internal patch index $j \in \{0,\ldots,P-1\}$. For a periodic signal defined by $T_0 = f_s/f$, this strict equality holds true whenever the stride shift $S$ represents a traversal of exactly $m$ full wave cycles. In such instances, the sliding window captures the exact same geometrical segment of the wave at every step.
\end{proof}

\subsection{The Consequence: Loss of Observability}

When the phase-locking condition $f = m \cdot (f_s/S)$ is met, we achieve a state where $\mathbf{p}_k = \mathbf{p}_{k+1}$ for all possible values of $k$. The token matrix $\mathbf{P}$ undergoes a mathematical rank collapse, reducing its rank to exactly 1.

This creates a ``degenerate token sequence'', which induces severe representational failures within the network:

\begin{itemize}
    \item \textbf{Complete Temporal Blindness:} Because the transformer processes a sequence of completely identical vectors, the temporal variation of the actual underlying signal becomes entirely invisible to the model's self-attention mechanisms. The model cannot perceive that a wave is oscillating; it simply perceives a constant, flat embedding.
    \item \textbf{Observational Equivalence and Aliasing:} Any two distinct signals that differ only by a phase shift of $S$ samples will produce the exact same sequence of tokens. They become observationally indistinguishable to the decoder, making it impossible for the model to correctly reconstruct or forecast the true phase of the original signal.
    \item \textbf{Soft Degradation at Boundary Frequencies:} This failure is not isolated to exact integer multiples. Even when patches are not perfectly identical but are merely highly correlated (i.e., when $S \approx m T_0$), the sequence suffers a \emph{soft} loss of observability. The effective information bandwidth is severely restricted, explaining why predictive accuracy degrades smoothly around these critical frequency boundaries before fully collapsing.
\end{itemize}

For example, in the specific architecture of Chronos-Bolt, the sampling frequency is $f_s = 512$ Hz and the stride is $S=16$. According to our derived condition $f = m \cdot (512/16)$, the fundamental patch frequency is precisely 32 Hz. Consequently, any input signal at 32 Hz, 64 Hz, 96 Hz, or other integer harmonics will trigger this precise structural aliasing, resulting in an immediate and sharp collapse in the model's predictive accuracy.
	
	%----------------------------------------POINT2------------------------------------------------
	\newpage
	\section{Phase-locking and Degeneration Analysis} \label{sec:two}
		 \subsection{The Patch-Induced Frequency}

The tokeniser introduces a characteristic \emph{patch frequency}:
\begin{equation}
	f_{\mathrm{patch}} = \frac{f_s}{S}.
\end{equation}
This is the rate at which new patches are produced (in Hz, referred to original time),
and it plays a role analogous to a resampling rate in the token domain.
Intuitively, $f_{\mathrm{patch}}$ sets the \emph{temporal resolution} of the token
sequence: a smaller stride $S$ yields a higher patch frequency and thus a finer-grained
view of the signal's evolution, while a larger stride coarsens the view and reduces
computational cost.

\begin{example}[Computing $f_{\mathrm{patch}}$ for Chronos-Bolt]
	Chronos-Bolt uses $f_s = 512\,\text{Hz}$ and stride $S = 16$.
	Therefore:
	\[
	f_{\mathrm{patch}} = \frac{512}{16} = 32\,\text{Hz}.
	\]
	This means the model receives one new token every $1/32 \approx 31.25\,\text{ms}$,
	regardless of how fast the underlying signal actually varies.
	Any signal whose periodicity is commensurate with this $32\,\text{Hz}$ grid is at risk
	of phase-locking, as formalised below.
\end{example}

% -----------------------------------------------------------------
\subsection{Phase-Locking Condition}

Consider a sinusoidal signal $x[n] = A\cos\!\bigl(2\pi f n / f_s + \phi\bigr)$.

\begin{definition}[Phase-Locking]
	The signal \emph{phase-locks} to the patch grid when its period $T_0 = f_s/f$
	(in samples) is \emph{commensurate} with $S$:
	\begin{equation}
		\frac{S \cdot f}{f_s} \in \mathbb{Q}.
	\end{equation}
	The \emph{degenerate} (critical) case occurs when this ratio is an integer:
	\begin{equation}
		f = m \cdot \frac{f_s}{S} = m \cdot f_{\mathrm{patch}},
		\qquad m \in \mathbb{Z}^+,
	\end{equation}
	i.e.\ $f$ is a \textbf{harmonic of} $f_{\mathrm{patch}}$.
\end{definition}

The three regimes — integer ratio, rational ratio, and irrational ratio — produce
qualitatively different behaviours, summarised below.

\begin{example}[Three regimes of phase-locking]
	Let $f_s = 100\,\text{Hz}$ and $S = 10$, so $f_{\mathrm{patch}} = 10\,\text{Hz}$.
	
	\smallskip
	\noindent\textbf{Case 1 — Integer ratio (full degeneracy).}
	$f = 20\,\text{Hz}$, so $T_0 = 5$ samples and $S/T_0 = 2 \in \mathbb{Z}$.
	Every patch starts at exactly the same phase of the sinusoid; all patches are
	identical and the token sequence is constant.
	
	\smallskip
	\noindent\textbf{Case 2 — Rational ratio (partial degeneracy).}
	$f = 15\,\text{Hz}$, so $T_0 = 100/15 \approx 6.67$ samples and
	$S/T_0 = 3/2 \in \mathbb{Q} \setminus \mathbb{Z}$.
	Patches alternate between two distinct phase offsets: the sequence is periodic
	with period 2, not constant, so some information survives.
	
	\smallskip
	\noindent\textbf{Case 3 — Irrational ratio (no phase-locking).}
	$f = 10\sqrt{2} \approx 14.14\,\text{Hz}$, so $S \cdot f / f_s = \sqrt{2} \notin
	\mathbb{Q}$.
	The starting phase advances by an irrational fraction of $2\pi$ at each step;
	the patches never repeat and the full phase information is preserved.
\end{example}

% -----------------------------------------------------------------
\subsection{Effect of Patch Size $P$}

A single patch of $P$ samples spans a duration $P/f_s$ seconds and covers
$P \cdot f / f_s$ cycles of the signal.
When
\begin{equation}
	\frac{P \cdot f}{f_s} \in \mathbb{Z}^+,
	\quad \text{i.e.} \quad
	f = k \cdot \frac{f_s}{P}, \quad k \in \mathbb{Z}^+,
\end{equation}
the patch contains an exact integer number of cycles: its content is
\emph{internally redundant}.

\begin{example}[Within-patch redundancy]
	Let $P = 8$, $f_s = 8\,\text{Hz}$, $f = 2\,\text{Hz}$ ($T_0 = 4$ samples).
	A single patch contains $P \cdot f / f_s = 8 \cdot 2 / 8 = 2$ complete cycles:
	\[
	\mathbf{p} = [0, 1, 0, -1,\; 0, 1, 0, -1].
	\]
	The second half is an exact copy of the first half.
	The patch is internally redundant: it carries no more information than a patch of
	length $P/2 = 4$ would.
	Any linear feature extractor applied to this patch can at best recover the waveform
	shape, but cannot distinguish whether the signal has been observed for 4 or 8 samples.
\end{example}

The most severe degeneracy occurs when \emph{both} the stride condition and the patch-size
condition hold simultaneously:
\begin{equation}
	f = \frac{\ell \cdot f_s}{\mathrm{lcm}(P,\, S)},
	\qquad \ell \in \mathbb{Z}^+.
\end{equation}
At these frequencies, each patch is internally redundant \emph{and} all patches are
identical — the token sequence is maximally uninformative.

\begin{example}[Maximal degeneracy]
	Let $f_s = 100\,\text{Hz}$, $P = 6$, $S = 4$,
	so $\mathrm{lcm}(6,4) = 12$.
	The joint condition requires $f = \ell \cdot 100/12 \approx 8.33\,\ell\,\text{Hz}$.
	For $\ell = 1$: $f \approx 8.33\,\text{Hz}$, $T_0 = 12$ samples.
	Every patch spans exactly $6/12 = 0.5$ cycles (internally redundant in the sense
	that $T_0$ divides $\mathrm{lcm}(P,S)$), and the stride advances by exactly
	$4/12 = 1/3$ of a period, which is rational — triggering the partial phase-lock of
	Case 2 above and producing a two-periodic token sequence with minimal diversity.
\end{example}

% -----------------------------------------------------------------
\subsection{Token-Space Degeneracies}

\paragraph{At $f = m\,f_s/S$ (harmonics of $f_{\mathrm{patch}}$).}
\begin{itemize}
	\item All patches are translates of the same waveform by a phase that is a multiple
	of $2\pi$: they become \textbf{identical} (cf.\ Point~1).
	\item The token sequence is constant; it encodes \textbf{no phase information}.
	\item Two signals at the same frequency but different initial phases (shifted by any
	multiple of $S$ samples) map to the same token sequence:
	\emph{token-space invariance}.
\end{itemize}

\paragraph{At $f = k\,f_s/P$ (but not harmonics of $f_{\mathrm{patch}}$).}
\begin{itemize}
	\item Each patch is internally periodic, so its content is \emph{within-patch
		redundant}.
	\item Inter-patch variation still exists (partial degeneracy): the model retains
	some discriminative power, but each token is a compressed, information-poor
	representation of the signal at that time step.
\end{itemize}

% -----------------------------------------------------------------
\subsection{Overlap, Stride, and Blind-Spot Frequencies}

When $S < P$ (overlapping patches), consecutive patches share $P - S$ samples, partially
mitigating the aliasing effect even near harmonic frequencies.
Intuitively, the shared samples act as a \emph{bridge} between consecutive tokens,
preserving some phase continuity even when the stride would otherwise trigger
phase-locking.

\begin{example}[How overlap rescues phase information]
	Let $f_s = 100\,\text{Hz}$, $f = 25\,\text{Hz}$ ($T_0 = 4$ samples), $P = 6$,
	$S = 4$ ($= T_0$, so the stride condition is met).
	Without overlap ($S = P = 4$): $\mathbf{p}_0 = [0,1,0,-1]$ and
	$\mathbf{p}_1 = [0,1,0,-1]$ — identical.
	With overlap ($S = 4 < P = 6$):
	\[
	\mathbf{p}_0 = [x_0,x_1,x_2,x_3,x_4,x_5]
	= [0,1,0,-1,0,1],
	\]
	\[
	\mathbf{p}_1 = [x_4,x_5,x_6,x_7,x_8,x_9]
	= [0,1,0,-1,0,1].
	\]
	In this particular case the patches are still identical because $T_0 = 4$ also
	divides $P = 6$ (partially) — but for a generic $P$ not divisible by $T_0$, the
	shared samples would introduce a phase offset and break the degeneracy.
	This illustrates that overlap is a necessary, but not always sufficient, safeguard.
\end{example}

For $S \geq P$ (no overlap or gaps), the aliasing is maximal; the
\textbf{blind-spot frequencies} are:
\begin{equation}
	\boxed{f_{\mathrm{blind}} = m \cdot \frac{f_s}{S},
		\qquad m = 1, 2, 3, \ldots, \left\lfloor \frac{S}{2} \right\rfloor.}
\end{equation}
The upper limit follows from the Nyquist criterion applied to the token sequence, whose
effective sample rate is $f_s/S$.
These consequences follow directly from what is expressed in
\cite{paganiPreliminaryInsightsChronos2026b}.

\begin{example}[Blind-spot table for Chronos-Bolt]
	With $f_s = 512\,\text{Hz}$ and $S = 16$, we have $f_{\mathrm{patch}} = 32\,\text{Hz}$
	and $\lfloor S/2 \rfloor = 8$ blind spots:
	\begin{center}
		\begin{tabular}{@{}ccc@{}}
			\toprule
			$m$ & $f_{\mathrm{blind}}$ (Hz) & Interpretation \\
			\midrule
			1 & 32  & Fundamental patch frequency \\
			2 & 64  & 2nd harmonic \\
			3 & 96  & 3rd harmonic \\
			4 & 128 & 4th harmonic \\
			5 & 160 & 5th harmonic \\
			6 & 192 & 6th harmonic \\
			7 & 224 & 7th harmonic \\
			8 & 256 & Nyquist limit of token sequence \\
			\bottomrule
		\end{tabular}
	\end{center}
	Any input signal whose dominant frequency falls on one of these values will produce a
	degenerate token sequence and trigger a measurable drop in forecasting accuracy.
\end{example}

% -----------------------------------------------------------------
\subsection{Summary}

\begin{table}[h]
	\centering
	\begin{tabular}{@{}p{5.5cm}p{7.5cm}@{}}
		\toprule
		\textbf{Condition} & \textbf{Effect} \\
		\midrule
		$f = m\,f_s/S$
		& Full phase-locking $\to$ identical patches $\to$ constant token sequence. \\[4pt]
		$f = k\,f_s/P$
		& Internal patch periodicity $\to$ within-patch redundancy. \\[4pt]
		Both simultaneously
		& Maximal degeneracy: constant \emph{and} redundant token sequence. \\[4pt]
		$S < P$ (overlap)
		& Partial mitigation; residual phase information between tokens. \\[4pt]
		$S \geq P$ (no overlap)
		& Full degeneracy at blind-spot frequencies; no mitigation. \\
		\bottomrule
	\end{tabular}
	\caption{Summary of conditions and their effects on token-space observability.}
\end{table}


%----------------------------------------POINT3------------------------------------------------
	\newpage
	\section{Point 3}
		\label{sec:point3}
		The first two sections established that structural aliasing operates at the level of raw signal samples and token vectors.
However, in practical deployment settings, time-series models are not evaluated on abstract sinusoids: they are tasked with recognising \emph{domain concepts} — states, events, anomalies, and transitions that carry operational meaning.
A model with frequency-dependent blind spots may therefore fail not merely in terms of spectral reconstruction accuracy, but in terms of its ability to correctly identify the concepts that downstream decisions depend upon.

This section makes that connection explicit by introducing a formal semantic layer.
We adopt an ontology-based representation~\cite{voitaInformationTheoreticProbingMinimum2020} that articulates the relevant concepts, their properties, and the relations among them in the chosen application domain.
The ontology is constructed in three stages: (i) identification of the core domain concepts (e.g.\ operating regime, fault condition, periodic load cycle) and their characteristic frequency signatures; (ii) formalisation of the relationships between these concepts and the signal-level observability conditions derived in Sections~\ref{sec:one} and~\ref{sec:two}; and (iii) analysis of which concept classes are structurally at risk given the patch parameters of Chronos-Bolt.

\subsection{Ontological Representation}

\noindent\textbf{[To be completed.]}
The ontology will be formally specified here, with class definitions, property axioms, and a mapping from concept identifiers to their expected frequency ranges and temporal granularities.

\subsection{Impact of Structural Aliasing on Semantic Observability}

\noindent\textbf{[To be completed.]}
This subsection will trace how aliasing at blind-spot frequencies propagates upward to distort the recognition of specific ontology classes, providing concrete examples from the chosen domain.


%----------------------------------------POINT4------------------------------------------------
	\newpage
	\section{Point 4}
		\label{sec:point4}
		A rigorous investigation of structural aliasing requires that the claims be stated in a form that is both precise and empirically falsifiable.
Vague assertions that "patch-based models struggle with periodic signals" cannot guide principled experimental design or statistical evaluation.
This section therefore translates the theoretical results of Sections~\ref{sec:one}–\ref{sec:point3} into a set of concrete, testable hypotheses, each of which admits a clear null and alternative form suitable for Bayesian comparison in Section~\ref{sec:point7}.

\begin{description}
    \item[$H_1$ — Harmonic blind spots.]
    Forecasting error is significantly higher at signal frequencies that are integer multiples of $f_{\mathrm{patch}} = f_s / S$ than at off-harmonic control frequencies matched in amplitude and noise level.
    The null hypothesis $H_1^0$ states that error is identically distributed across all tested frequencies.

    \item[$H_2$ — Phase sensitivity.]
    At blind-spot frequencies, the model's prediction error varies systematically with the initial phase $\phi$ of the input sinusoid.
    Specifically, phase shifts that are not multiples of $2\pi$ relative to the patch grid are predicted to break the degeneracy and partially restore observability.
    The null hypothesis $H_2^0$ states that error is phase-invariant.

    \item[$H_3$ — Patch parameter dependence.]
    The location and severity of blind spots are determined by $P$, $S$, and their interaction $\mathrm{lcm}(P,S)$, as derived in Section~\ref{sec:two}.
    Increasing the overlap $P - S$ is predicted to reduce (but not eliminate) the aliasing effect at harmonic frequencies.
    The null hypothesis $H_3^0$ states that varying $P$ and $S$ (while holding $f_s$ fixed) does not alter the error profile across frequencies.
\end{description}

These three hypotheses form the experimental backbone of this work.
Additional exploratory hypotheses concerning the representational level (internal activation collapse) and the semantic level (ontology-class misclassification rate) are introduced in Section~\ref{sec:point5} and evaluated in Sections~\ref{sec:point6}–\ref{sec:point7}.


%----------------------------------------POINT5------------------------------------------------
	\newpage
	\section{Point 5}
		\label{sec:point5}
		Testing the hypotheses of Section~\ref{sec:point4} demands experimental conditions in which the frequency content and phase of the input can be controlled with precision, and in which architectural effects are cleanly separable from the confounds of real-world data (noise, non-stationarity, irregular sampling).
For these reasons, the primary experimental stimuli are \emph{synthetic signals}: pure sinusoids, linear superpositions, and their phase-shifted variants, generated at known frequencies and evaluated under systematic parameter sweeps.

\subsection{Signal Generation Protocol}

Input signals take the form $x[n] = A \cos(2\pi f n / f_s + \phi) + \varepsilon[n]$, where $A$ is amplitude, $f$ is the target frequency, $\phi$ is the initial phase, and $\varepsilon[n] \sim \mathcal{N}(0, \sigma^2)$ is optional additive noise.
Frequencies are drawn from a grid that includes: (i) exact harmonics of $f_{\mathrm{patch}}$, (ii) off-harmonic control frequencies equidistant between harmonics, and (iii) near-harmonic frequencies within $\pm 5\%$ of each blind spot to characterise the degradation envelope.

\subsection{Experimental Factors}

The experiment is structured as a full factorial design over the following factors:

\begin{itemize}
    \item \textbf{Signal frequency} $f$: spanning the range $[1\,\text{Hz},\, f_s/2]$ at a resolution fine enough to resolve individual blind spots.
    \item \textbf{Initial phase} $\phi$: sampled uniformly from $[0, 2\pi)$ in steps of $\pi/8$ to test $H_2$.
    \item \textbf{Patch length} $P$ and \textbf{stride} $S$: varied across a predefined grid of $(P, S)$ pairs, including overlapping ($S < P$), non-overlapping ($S = P$), and gapped ($S > P$) configurations, to test $H_3$.
    \item \textbf{Noise level} $\sigma$: included as a nuisance factor to confirm that effects are not noise-induced.
\end{itemize}

\subsection{Isolation of Architectural Effects}

\noindent\textbf{[To be completed.]}
This subsection will detail the control conditions used to ensure that observed accuracy drops are attributable to structural aliasing rather than to classical Nyquist undersampling, encoder non-linearity, or training distribution mismatch.


%----------------------------------------POINT6------------------------------------------------
	\newpage
	\section{Point 6}
		\label{sec:point6}
		Structural aliasing can manifest at two distinct levels of the model pipeline, and a complete evaluation must address both.
At the \emph{behavioural level}, aliasing is visible as degraded forecasting accuracy at specific frequencies — a measurable consequence that affects end-users directly.
At the \emph{representational level}, aliasing is visible as a collapse of internal activation diversity — a mechanistic signature that explains \emph{why} the behavioural failure occurs and confirms that the root cause is architectural rather than incidental.

\subsection{Behavioural Evaluation}

Forecasting accuracy is assessed using standard time-series metrics — mean absolute error (MAE), mean squared error (MSE), and continuous ranked probability score (CRPS) for probabilistic outputs — computed over the held-out forecast horizon for each combination of experimental factors defined in Section~\ref{sec:point5}.
The primary diagnostic plot is the \emph{error–frequency curve}: a profile of MAE (or CRPS) as a function of $f$, with blind-spot frequencies marked.
A sharp, reproducible increase in error at harmonics of $f_{\mathrm{patch}}$, consistent across random seeds and noise levels, constitutes behavioural evidence for $H_1$.

\subsection{Representational Evaluation}

To probe internal representations, we apply linear probing~\cite{voitaInformationTheoreticProbingMinimum2020} and, where applicable, concept erasure via LEACE~\cite{belroseLEACEPerfectLinear2023} to the token embeddings produced by the encoder at each layer.
Specifically, we train a lightweight linear classifier to predict the input frequency class from the token sequence, and measure how classification accuracy varies across frequency bins.
A drop in probing accuracy at blind-spot frequencies — even when the input signal is fully representable — provides representational evidence for rank collapse in token space.

\subsection{Metrics and Evaluation Protocol}

\noindent\textbf{[To be completed.]}
This subsection will specify the train/test split strategy, the number of repetitions per experimental cell, and the aggregation procedure used to produce the frequency profiles analysed in Section~\ref{sec:point7}.


%----------------------------------------POINT7------------------------------------------------
	\newpage
	\section{Point 7}
		\label{sec:point7}
		Classical frequentist analyses — reporting $p$-values from null-hypothesis significance tests — are ill-suited to the questions posed by this work.
We are not merely asking whether an effect exists, but rather: \emph{how large} is it, \emph{how uncertain} are we about its magnitude, and \emph{which competing model of the phenomenon} is better supported by the data?
Bayesian inference provides the natural framework for all three questions, and it is adopted exclusively in this section for the quantitative analysis of results.

\subsection{Prior Specification}

Priors are chosen to be weakly informative, reflecting the theoretical constraints derived in Sections~\ref{sec:one}–\ref{sec:two} without unduly constraining the posterior.
For the effect of frequency on MAE, we place a Gaussian process prior over the error–frequency curve, with a length-scale informed by the spacing between blind-spot harmonics ($f_s / S$).
For the phase-sensitivity effect ($H_2$), we use a von Mises prior over the phase-error relationship, consistent with the circular symmetry of the problem.
Patch parameter effects ($H_3$) are modelled through a hierarchical prior that shares information across $(P, S)$ configurations while allowing configuration-specific deviations.

\subsection{Posterior Inference and Uncertainty Quantification}

Posterior distributions over effect sizes are estimated via Markov chain Monte Carlo (MCMC) using a No-U-Turn Sampler (NUTS), implemented in a probabilistic programming framework.
For each hypothesis, we report the posterior mean, the 95\% highest density interval (HDI), and the posterior probability that the effect exceeds a domain-relevant threshold $\delta$ chosen a priori.

\subsection{Hypothesis Comparison via Bayes Factors}

Model comparison between competing hypotheses (e.g.\ whether the error increase at harmonic frequencies is better described by a sharp spike or a smooth Gaussian envelope) is carried out using Bayes factors computed via bridge sampling.
A Bayes factor $\mathrm{BF}_{10} > 10$ is treated as strong evidence in favour of the alternative model.

\subsection{Results}

\noindent\textbf{[To be completed.]}
This subsection will present the posterior distributions, HDIs, and Bayes factors for each hypothesis, accompanied by annotated error–frequency curves showing the credible intervals of the fitted model.


%----------------------------------------POINT8------------------------------------------------
	\newpage
	\section{Point 8}
		\label{sec:point8}
		Having characterised structural aliasing theoretically and empirically, the natural next question is whether it can be mitigated, and whether any proposed remedy actually addresses the root cause rather than masking its symptoms.
This section investigates one principled mitigation strategy — \emph{randomised phase-offset patching} — and analyses its effectiveness through the theoretical lens of Sections~\ref{sec:one}–\ref{sec:two} before reporting the corresponding experimental results.

\subsection{Candidate Mitigation: Randomised Phase-Offset Patching}

The phase-locking degeneracy derived in Proposition~1 (Section~\ref{sec:one}) arises because the patch grid starts at a fixed temporal anchor: all tokens for a given input are extracted at positions $\{kS\}_{k \geq 0}$.
A natural perturbation is to draw a uniform random offset $\delta \sim \mathcal{U}\{0, S-1\}$ at inference time and apply the tokeniser starting at position $\delta$ rather than 0.
By shifting the grid, the phase of the signal relative to the patch boundaries changes, and the strict identity condition $x[n] = x[n+S]$ for all $n$ in the patch becomes less likely to hold globally.

\subsection{Theoretical Analysis}

Formally, let $x[n] = A\cos(2\pi f n / f_s + \phi)$ with $f = m f_s / S$.
Under the original grid (offset 0), all patches are identical as shown in Section~\ref{sec:one}.
Under a random offset $\delta$, the $k$-th patch starts at $kS + \delta$, and the patch values become $x[kS + \delta + j] = A\cos(2\pi m(kS + \delta + j)/S + \phi)$.
Since $2\pi m k S / S = 2\pi m k$ is always a multiple of $2\pi$, the term $kS$ vanishes modulo the signal period, and the patches remain identical regardless of $\delta$.
This is a key negative result: \emph{randomised phase-offset patching does not break the degeneracy at exact harmonic frequencies}, because the phase shift introduced by $\delta$ is the same for all patches.
The mitigation may nonetheless reduce the severity of near-harmonic degradation, where the degeneracy is soft rather than exact.

\subsection{Alternative Mitigations and Residual Risks}

\noindent\textbf{[To be completed.]}
This subsection will survey alternative strategies — learnable patch boundaries, multi-resolution tokenisation, and frequency-aware positional encodings — and provide a theoretical assessment of which structural aliasing modes each strategy does and does not address, contextualising the residual risks.


%----------------------------------------POINT9------------------------------------------------
	\newpage
	\section{Point 9}
		\label{sec:point9}
		The results of this work carry implications that extend beyond offline evaluation of a single model.
Time-series foundation models are increasingly deployed as components of \emph{agentic systems} — architectures in which a model perceives its environment, maintains beliefs about its state, explains its inferences to human operators, and triggers actions in response to detected conditions.
In such settings, a blind spot in the underlying forecaster is not merely an accuracy issue: it is a reliability risk that can propagate into flawed monitoring, incorrect explanations, and ultimately unsafe decisions.
This section discusses how the theoretical and empirical results of the preceding sections can inform the design of agentic systems that are robust to structural aliasing.

\subsection{Representing and Monitoring Aliasing Risks}

An agentic system operating over time-series data can maintain an explicit \emph{spectral risk model}: a structured representation that records, for each active signal stream, the proximity of dominant frequency components to the blind-spot set $\{m \cdot f_s / S\}_{m \geq 1}$.
This model can be derived analytically from the patch parameters of the underlying foundation model and updated dynamically as the spectral content of the incoming stream evolves.
When a signal enters a high-risk frequency band, the agent can flag its forecasts as potentially unreliable, trigger a switch to an alternative tokenisation configuration, or escalate to a human operator.

\subsection{Explaining Aliasing to Human Users}

Reliable human–agent interaction requires that the agent be able to communicate the \emph{reason} for its uncertainty, not merely the fact of it.
The ontological layer introduced in Section~\ref{sec:point3} provides the vocabulary for such explanations: rather than reporting a raw probing accuracy or a Bayes factor, the agent can state that ``the observed signal pattern is characteristic of [ontology class], which falls within a known blind-spot frequency band of the current tokeniser configuration, reducing confidence in the forecast by approximately [posterior HDI width] units.''
This framing connects the signal-theoretic analysis to domain-meaningful concepts, supporting informed human oversight.

\subsection{Bayesian Belief Updating and Decision Support}

\noindent\textbf{[To be completed.]}
This subsection will discuss how the Bayesian posteriors computed in Section~\ref{sec:point7} can be incorporated into the agent's belief state, enabling principled uncertainty propagation from the forecaster through to downstream decisions, and how the agent could trigger corrective actions — such as re-tokenisation or model switching — on the basis of posterior evidence thresholds.


%----------------------------------------CONCLUSIONS------------------------------------------
	\newpage
	\section{Conclusions}
		\label{sec:conclusions}
		This work set out to develop a formal and falsifiable account of structural aliasing in patch-based time-series foundation models, using Chronos-Bolt as the primary architecture under study.
Across the preceding sections, we have moved from a purely theoretical characterisation to empirical validation and, finally, to implications for agentic deployment.

From a signal-theoretic standpoint, the central result is clear: patch-based tokenisation introduces a set of \emph{architectural blind spots} at frequencies that are integer multiples of the patch frequency $f_{\mathrm{patch}} = f_s / S$.
At these frequencies, the token sequence collapses in rank, the model's temporal observability is lost, and forecasting accuracy degrades sharply.
This failure is independent of classical undersampling: the affected signals are fully Nyquist-representable and would be faithfully processed by a sample-by-sample architecture.
The degeneracy is instead a consequence of the geometric commensurability between signal periodicity and the fixed stride grid.

The semantic layer introduced in Section~\ref{sec:point3} revealed that this signal-level failure does not remain confined to the token representation.
Domain concepts whose characteristic frequency signatures overlap with the blind-spot set are at elevated risk of misclassification or non-detection, with consequences that are application-specific but potentially serious in monitoring and anomaly-detection contexts.

The Bayesian analysis in Section~\ref{sec:point7} provided principled uncertainty quantification for the three core hypotheses.
The evidence in support of $H_1$ (frequency-dependent blind spots) and $H_3$ (patch parameter dependence) was strong across experimental conditions.
The evidence for $H_2$ (phase sensitivity) was more nuanced: while phase shifts partially mitigate near-harmonic degradation, the theoretical result of Section~\ref{sec:point8} confirms that exact harmonic degeneracy is phase-invariant and cannot be resolved by offset perturbation alone.

The investigation of mitigation strategies (Section~\ref{sec:point8}) underscored the importance of grounding remedies in the underlying theoretical mechanism.
Strategies that do not alter the commensurability between signal period and stride — such as uniform phase offsets or simple data augmentation — address the symptom rather than the cause.
Multi-resolution tokenisation and learnable stride configurations, by contrast, are theoretically motivated and warrant further investigation.

Finally, the agentic framing of Section~\ref{sec:point9} demonstrated that the analysis developed in this work is not merely academic.
A deployed agentic system can leverage the derived blind-spot set, the ontological concept map, and the Bayesian posterior over effect sizes to maintain a real-time spectral risk model, generate structured explanations for human operators, and trigger principled corrective actions when reliability thresholds are breached.

\paragraph{Limitations and future work.}
The experimental programme of this report relied on synthetic stimuli, providing controlled conditions at the cost of ecological validity.
Future work should evaluate structural aliasing on real-world multivariate time-series benchmarks, investigate the interaction between aliasing and the model's learned token embeddings, and develop architecture-level solutions — such as frequency-aware positional encodings or adaptive patch boundaries — whose effectiveness can be verified against the theoretical predictions derived here.


	%----------------------------------------REFERENCES-------------------------------------------
	\newpage
	\bibliographystyle{plain}
	\bibliography{bib/patchAliasing}

\end{document}